\centerline{\bf \Huge Оценка + пример I}

\begin{flushright}
        \scriptsize{Никогда не думайте, что вы уже всё знаете. И как бы высоко ни оценивали вас, всегда имейте мужество сказать себе: «Я невежда».}
\end{flushright}

\problem{\textbf{1}} Какое наибольшее число трёхклеточных уголков можно вырезать из клетчатого квадрата 10$\times$10?

\problem{\textbf{2}} Какое наибольшее количество ладей можно поставить на шахматную доску так, чтобы они не били друг друга? (ладья бьет по вертикали и горизонтали)

\problem{\textbf{3}} Найдите наименьшее натуральное число кратное 10, сумма цифр которого равна 10

\problem{\textbf{4}} В пруд пустили 30 щук, которые стали кушать друг друга. Щука считается сытой, если она съела хотя бы трёх щук. Какое наибольшее количество щук могло насытиться, если съеденные сытые щуки при подсчёте тоже учитываются?

\problem{\textbf{5}} Найдите наибольшее число, любые две последовательные цифры которого образуют точный квадрат в том порядке, в котором они стоят. (Точный квадраты это: 0, 1, 4, 9, 16, 25...)

\problem{\textbf{6}} На клетчатой доске 100×100 закрасили Х прямоугольников 1×2.
Оказалось, что в каждой строке и в каждом столбце есть хотя
бы одна закрашенная клетка. При каком наименьшем Х это
возможно?

\problem{\textbf{7}} Семизначное число называется неразложимым, если оно не
раскладывается в произведение двух четырёхзначных чисел. Какое
наибольшее количество неразложимых семизначных чисел может идти
подряд?
